% Values: PRECISE, FAITHFUL, TRACTABLE, ILLUMINATING, ECONOMICAL
% Evidence planned:
% - PRECISE: a rule for every production of the bare module grammar
%   (Figures fig:grammar-modules and fig:grammar-top-level-declarations):
%   19 rules in 9 groups; the exposing list of a header is covered by the
%   operation Expose, every other production by a rule; interfaces, provider
%   maps, ambiguity and the Port predicate defined before the first rule.
% - FAITHFUL: default imports first; replacement of explicitly exposed types,
%   merging of everything else, lazy ambiguity, no qualified
%   self-reference (A28, A38); open, closed and private types, record alias
%   constructors of status v, `Alias(..)' and constructors by name rejected;
%   private types in exposed types legal (A29); ports as Elm 0.19.1 with
%   Bytes added (D11, A92), with the table of what can and cannot be sent; main left to
%   Chapter 9 (A32); port names program-wide left to Chapter 9 (A84).
% - TRACTABLE: the property that a module generates only new type names and
%   that the type environment of its type declarations is determined up to the
%   choice of those names; interfaces made by Expose are well formed.
% - ILLUMINATING: comments on the fixed point of the type declarations, the
%   two port rules, replacement versus merging, ambiguity, the record alias
%   constructor, and the absence of qualified self-reference.
% - ECONOMICAL: no signatures, realisation or enrichment (D6); one rule for a
%   header with and without `port'; options merge `as' and `exposing'.

\chapter{Static Semantics for Modules}\label{ch:static-modules}

This chapter defines the static semantics of Modules. % D6
It introduces interfaces, which record what a module makes available to the
modules that import it, module environments, port environments and static
bases, together with provider maps, which record which imported module
supplies each identifier that a module may use. The rules infer that a module
elaborates, in a basis, to the part of a basis it contributes and to the
environment of its own declarations; on the way they infer what the imports,
exposing lists, top-level declarations and type declarations of the module
elaborate to. Since Elm modules export top-level names only, there are no
signatures, functors or sharing, and no matching of an environment against a
specification. The admissible types of ports are defined in
Section~\ref{sec:static-modules-ports}. % D11

\section{Semantic Objects}\label{sec:static-modules-objects}

The simple objects for the static semantics of Modules are exactly as for the
Core (Section~\ref{sec:static-core-simple}). The compound objects are those
for the Core, augmented by those in Figure~\ref{fig:static-modules-compound}.

\begin{figure}[ht]
\centering
\begin{objtable}
  \objdef{I}{\Interface}{\Env}
  \objdef{\ME}{\ModEnv}{\cls{ModId} \finmap \Interface}
  \objdef{\PE}{\PortEnv}{\cls{Var} \finmap \{\smdirout, \smdirin\} \times \Type}
  \objdef{B\ \text{or}\ T, \ME, \PE}{\Basis}{\Fin(\TyName) \times \ModEnv \times \PortEnv}
  \objdef{\pi}{\smProvMap}{\smLongId \finmap (\Fin(\cls{ModId}) \setminus \{\emptyset\})}
\end{objtable}
\caption{Further Compound Semantic Objects}\label{fig:static-modules-compound}
\end{figure}

An \emph{interface} $I$ is an environment whose qualified environment is empty.
It is what a module \emph{exposes}: the types, values and constructors that
another module may name after importing it. Every type that a module declares
is either \emph{private}, if $\tycon \notin \Dom(\TE \of I)$, or exposed. An
exposed custom type is \emph{open} if its constructors belong to $\Dom(\VE \of
I)$, and \emph{closed} if none of them does; in both cases $I(\tycon)$ is the
type structure of the type, constructors included. An exposed type alias is
bound in $\TE \of I$, and an exposed record alias brings its constructor, of
status $\isv$, into $\VE \of I$. In a type structure $(\theta, \VE)$ of an
interface, $\VE$ contains a constructor of status $\isc$ exactly when the type
constructor is a custom type, since a custom type has at least one constructor
and an alias has none of status $\isc$
(Section~\ref{sec:static-modules-rules}).
% DECIDED: a closed type keeps its constructors in its type structure, but they
% are absent from VE of I, so no importer can name them. This keeps Interface =
% Env (registry) and tells a closed custom type from an alias, which `T(..)' in
% an import needs (BAD IMPORT for an alias).

A \emph{module environment} $\ME$ maps the identifier of each module elaborated
so far to its interface. A \emph{port environment} $\PE$ maps the name of each
port to its \emph{direction}, $\smdirout$ for an outgoing port and $\smdirin$
for an incoming one, and to its \emph{payload type}, the type of the values it
carries. A \emph{static basis} $B = T, \ME, \PE$ consists of the set of type
names generated so far, the interfaces of the modules elaborated so far, and
the ports they declare; port names are global to a program
(Chapter~\ref{ch:programs}). % A84

We write $\smId$ for the union of $\cls{Var} \cup \cls{Con}$ and
$\cls{TyCon}$, and $\smLongId$ for the union of $\cls{LongVar} \cup
\cls{LongCon}$ and $\cls{LongTyCon}$; as in the Core, the unions are disjoint,
so a type constructor and a value constructor spelt alike are different members
of $\smId$. We let $\smid$ range over $\smId$ and $\smlongid$ over $\smLongId$,
and $N$ over $\Fin(\smLongId)$. For an environment $E$ whose qualified
environment is empty, $\Dom E$ means $\Dom(\TE \of E) \cup \Dom(\VE \of E)$, a
subset of $\smId$. For a set $N \subseteq \smId$ and a module identifier
$\modid$, $\modid.N$ means $\{\modid.\smid \,;\, \smid \in N\}$.

A \emph{provider map} $\pi$ maps each identifier that the imports of a module
make available, unqualified or qualified, to the set of modules that provide
it. An identifier with more than one provider is \emph{ambiguous}.

The operations of projection, injection and modification are as for the Core
(Section~\ref{sec:static-core-projection}). Moreover, for an environment $E$ we
write $C \oplus E$ for $C + (\tynames E, \emptyset, E)$, that is, $E \of C$
modified by $E$ and $T \of C$ extended with the type names of $E$; this
extends the notation $C \oplus \TE$ of the Core.

\section{Interfaces}\label{sec:static-modules-interfaces}

The interface that a module exposes is determined by the environment $E$ of
its own declarations and by the exposing list of its header. For an
environment $E$ whose qualified environment is empty, $\Expose(E, \explist)$
is defined as follows:
\[
\begin{array}{@{}l@{\ =\ }l@{}}
  \Expose(E, \kw{(..)}) & E\\
  \Expose(E, \kw{(}\expitem_1\kw{,}\ \cdots\kw{,}\ \expitem_n\kw{)})
    & \Expose(E, \expitem_1) + \cdots\\
  \multicolumn{1}{@{}l@{}}{} & \quad{} + \Expose(E, \expitem_n)\\
  \Expose(E, \var) & \{\var \mapsto E(\var)\}\ \text{in}\ \Env\\
  \Expose(E, \tycon) & (\emptymap,\ \{\tycon \mapsto (\theta, \VE)\},\ \VE_{\isv})\\
  \Expose(E, \tycon\,\kw{(..)}) & (\emptymap,\ \{\tycon \mapsto (\theta, \VE)\},\ \VE)
\end{array}
\]
where, in the last two lines, $E(\tycon) = (\theta, \VE)$ and $\VE_{\isv} =
\{\con \mapsto (\sigma, \isv) \,;\, \VE(\con) = (\sigma, \isv)\}$. It is defined when
every identifier of an item is bound in $E$, which the syntactic restrictions
of a header ensure (Section~\ref{sec:syntax-modules-restrictions}); a header
exposes no operator. Thus \kw{(..)} exposes every type open and every value,
constructor and port; an item $\tycon$ exposes a custom type closed and a
record alias with its constructor; and an item $\tycon\,\kw{(..)}$ exposes a
custom type open. No condition relates the exposed values to the exposed
types: a value may be exposed whose type mentions a private type. % A29

An interface $I$ is \emph{well formed} if, for every $(\theta, \VE) \in
\Ran(\TE \of I)$ and every $\con$ with $\VE(\con) = (\sigma, \isv)$, also $(\VE
\of I)(\con) = (\sigma, \isv)$. The rules rely on the following property,
which may easily be demonstrated from the rules for type declarations.

\begin{property}
If $E$ is the environment to which the top-level declarations of a module
elaborate, and $I = \Expose(E, \explist)$ is defined, then $I$ is well formed
and $I \subseteq E$ componentwise.
\end{property}

\noindent The interfaces of the initial basis
(Appendix~\ref{ch:appendix-static-basis}) are well formed.

\section{Provider Maps}\label{sec:static-modules-providers}

The imports of a module are read in order, and each one adds to a provider map.
For a module identifier $\modid$ and sets $N, N' \subseteq \smLongId$, the
provider map $\pi \smmerge_{\modid} (N, N')$ is $\pi$ in which $\modid$
\emph{replaces} the providers of the identifiers in $N$ and \emph{joins} the
providers of those in $N'$: if $\pi' = \pi \smmerge_{\modid} (N, N')$, then
\[
  \pi'(\smlongid) =
  \begin{cases}
    \{\modid\} & \text{if } \smlongid \in N,\\
    \pi(\smlongid) \cup \{\modid\} & \text{if } \smlongid \in N' \setminus N,\\
    \pi(\smlongid) & \text{if } \smlongid \in \Dom\pi \setminus (N \cup N'),
  \end{cases}
\]
where $\pi(\smlongid)$ is taken to be $\emptyset$ when $\smlongid \notin
\Dom\pi$. % A28

For $\modid \in \Dom\ME$, let $\TE_{\modid} = \TE \of \ME(\modid)$ and
$\VE_{\modid} = \VE \of \ME(\modid)$, except that for the module \kw{List}
of the initial basis
\[
\begin{array}{@{}l@{\ =\ }l@{}}
  \TE_{\kw{List}} & \TE_0 + \TE \of \ME(\kw{List}),\\
  \VE_{\kw{List}} & \VE_0 + \VE \of \ME(\kw{List}),
\end{array}
\]
where $E_0 = (\emptymap, \TE_0, \VE_0)$ is the environment of the type
constructor \kw{List} and the list constructors \kw{[]} and \kw{::} given in
Appendix~\ref{ch:appendix-static-basis}. % audit X10: E0 by reference
The environment that a
provider map $\pi$ gives in $\ME$ is $\smenv(\ME, \pi) = (\QE, \TE, \VE)$, where
\[
\begin{array}{@{}l@{\ =\ }l@{}}
  \TE & \{\tycon \mapsto \TE_{\modid}(\tycon) \,;\, \pi(\tycon) = \{\modid\}\}\\
  \VE & \{\smid \mapsto \VE_{\modid}(\smid) \,;\, \smid \in \cls{Var} \cup \cls{Con},\
        \pi(\smid) = \{\modid\}\}\\
  \QE & \{\modid' \mapsto (\emptymap, \TE_{\modid'}', \VE_{\modid'}') \,;\,
        \modid'.\smid \in \Dom\pi \text{ for some } \smid\}
\end{array}
\]
and, for each such $\modid'$,
\[
\begin{array}{@{}l@{\ =\ }l@{}}
  \TE_{\modid'}' & \{\tycon \mapsto \TE_{\modid}(\tycon) \,;\,
        \pi(\modid'.\tycon) = \{\modid\}\}\\
  \VE_{\modid'}' & \{\smid \mapsto \VE_{\modid}(\smid) \,;\,
        \smid \in \cls{Var} \cup \cls{Con},\ \pi(\modid'.\smid) = \{\modid\}\}.
\end{array}
\]
Thus an ambiguous identifier is bound in no component of $\smenv(\ME, \pi)$, and
a phrase that uses it does not elaborate, while one that does not use it is
unaffected. % A38
The rules add a module $\modid$ to $\pi(\smlongid)$ only when the last part of
$\smlongid$ is bound by $\TE_{\modid}$ or $\VE_{\modid}$, so $\smenv(\ME, \pi)$
is defined for every provider map that arises.

Before any import is read, the type \kw{List} and the list constructors of
$E_0$ are available, provided by the module \kw{List}:
\[
  \pi_0 = \{\kw{List} \mapsto \{\kw{List}\},\ \kw{[]} \mapsto \{\kw{List}\},\
  \kw{::} \mapsto \{\kw{List}\}\},
\]
where the first \kw{List} is a type constructor and the others are module
identifiers.
% Survey "Scoping and Recursion": the type List is in scope before any import,
% as an entry from the module List; static-core: the initial basis binds [] and
% :: with status c.

Operators are not bound in environments: the type scheme and status of an
operator are given by the static operator table $\OTstat$ of the initial basis,
in every module alike (Section~\ref{sec:static-core-projection}). For a module
identifier $\modid$, $\smops(\modid) \subseteq \Dom\OTstat$ is the set of
operators that the module $\modid$ declares. For a module of the initial basis
it is given in Appendix~\ref{ch:appendix-static-basis}
(Figure~\ref{fig:static-basis-operators}); for every other module it is empty,
since no program declares an operator. % A16; audit X1

The default imports of Figure~\ref{fig:default-imports}, taken in order, form a
phrase of the class $\cls{Imports}$, which we write $\imports_0$. For any
$\imports$, we write $\imports_0\ \imports$ for the import list consisting of
the default imports followed by those of $\imports$.

\section{Port Types}\label{sec:static-modules-ports}

The values that may cross a port are those of the types $\tau$ for which
$\Port(\tau)$ holds. % D11
$\Port$ is the smallest predicate on types that satisfies the conditions of
Figure~\ref{fig:static-modules-port}; it holds for no other type. In the
figure, $\tMaybe$, $\tArray$, $\tBytes$ and $\tValue$ are the type names of
the initial basis for \kw{Maybe}, \kw{Array}, \kw{Bytes.Bytes} and
\kw{Json.Encode.Value}
(Appendix~\ref{ch:appendix-static-basis}), whatever a program binds to type
constructors of the same spelling. The type \kw{Json.Decode.Value} is an alias
of \kw{Json.Encode.Value}, and so denotes $\tValue$ as well.

\begin{figure}[ht]
\centering
\[
\begin{array}{@{}l@{\qquad}l@{}}
  \Port(\tau) & \text{if } \tau \in \{\tInt, \tFloat, \tBool,\\
    & \phantom{\text{if } \tau \in \{} \tString, \tBytes, \tunit, \tValue\}\\[2pt]
  \Port(\tau\,\tList) & \text{if } \Port(\tau)\\[2pt]
  \Port(\tau\,\tMaybe) & \text{if } \Port(\tau)\\[2pt]
  \Port(\tau\,\tArray) & \text{if } \Port(\tau)\\[2pt]
  \Port((\tau_1, \tau_2)\tTuple{2}) & \text{if } \Port(\tau_1) \text{ and } \Port(\tau_2)\\[2pt]
  \Port((\tau_1, \tau_2, \tau_3)\tTuple{3})
    & \text{if } \Port(\tau_1),\ \Port(\tau_2)\\
    & \text{and } \Port(\tau_3)\\[2pt]
  \Port(\rec{\varrho}) & \text{if } \Port(\varrho(\lab))\\
    & \text{for every } \lab \in \Dom\varrho
\end{array}
\]
\caption{Port Types}\label{fig:static-modules-port}
\end{figure}

Types are semantic objects, in which every alias has been replaced by its
expansion (Section~\ref{sec:static-core-types}), so an alias of a port type is
a port type. Every port type is built from type names of the initial basis
alone. Figure~\ref{fig:static-modules-port-table} states the same predicate
by kinds of type, with the reason each excluded kind is excluded; the
encoding of the values of each port type as boundary values, which are JSON
values in which byte sequences may stand, is given in
Section~\ref{sec:effects-ports}. % A92

\begin{figure}[ht]
\centering
\small
\begin{tabular}{@{}p{0.36\textwidth}@{\quad}p{0.17\textwidth}@{\quad}p{0.37\textwidth}@{}}
\hline
Type & Port? & Reason\\
\hline
\kw{Int}, \kw{Float}, \kw{String}, \kw{Bool} & yes & JSON numbers, strings and Booleans\\
\kw{Bytes} & yes & byte sequences, which are not JSON values\\ % A92
\kw{()} & yes & JSON \kw{null}\\
\kw{Json.Encode.Value}, \kw{Json.Decode.Value} & yes & any JSON value, never
  a byte sequence\\
\kw{List} $t$, \kw{Array} $t$ & if $t$ is & JSON arrays\\
\kw{Maybe} $t$ & if $t$ is & \kw{null} or the value\\
2- and 3-tuples & if each component is & JSON arrays\\
closed records & if each field is & JSON objects\\
\hline
\kw{Char} & no & excluded in Standard Elm\\
function types & no & no JSON representation\\
type variables, constrained ones included & no & no JSON representation\\
extensible records & no & unknown fields\\
custom types other than \kw{Bool} and \kw{Maybe}: \kw{Result}, \kw{Order},
  \kw{Never}, \kw{Dict}, \kw{Set}, and every custom type of a program & no &
  no canonical encoding\\
\kw{Cmd}, \kw{Sub}, \kw{Task}, and other abstract types of the basis & no &
  no JSON representation\\
\hline
\end{tabular}
\caption{What can and cannot be sent over a port}\label{fig:static-modules-port-table}
\end{figure}
% A92 (superseding A18, A75): Bytes is a payload type, alone or inside any
% admissible type; implementations that reject it are recorded in Appendix E.
% Survey "Ports (Static)", payload types: Char, type variables, functions,
% extensible records and every custom type but Bool and Maybe are rejected.

\section{Inference Rules}\label{sec:static-modules-rules}

As for the Core, the rules of the static semantics of Modules allow sentences
of the form
\[
  \jdg{A}{\mathit{phrase}}{A'}
\]
to be inferred, where in this case $A$ is a basis, a basis paired with a
provider map, an environment, a module identifier paired with an interface, a
context, or a context paired with a type, and $A'$ is a semantic object or a
tuple of semantic objects. The sentence is pronounced ``$\mathit{phrase}$
elaborates to $A'$ in $A$''. Hypotheses not of this form are side conditions.
The convention for options is the same as for the Core static semantics.

The rules use the following notation for phrases. For a type-variable sequence
$\tyvarseq = \tyvar_1 \cdots \tyvar_k$ that parametrises a phrase (the
$\conbind$ of a custom-type binding or the $\ty$ of a type-alias binding), we
write $\tyvarseq = \tyrowvar^{(k)}$ when, for each $i$, $\tyrowvar_i$ is the
variable named by $\tyvar_i$ (Section~\ref{sec:static-core-simple}) that is a
row variable if $\tyvar_i$ occurs in the phrase as the extension variable of a
record type expression $\kw{\kwlb}\ \tyvar_i\ \kw{|}\ \tyrow\ \kw{\kwrb}$, and a
type variable otherwise. For a type constructor $\tycon$ we write $\con =
\tycon$ when $\con$ is spelt as $\tycon$. The partial function $\smlabs$ is
defined on a type expression $\ty$ that is a record type expression
$\kw{\kwlb}\ \lab_1\ \kw{:}\ \ty_1\kw{,}\ \cdots\kw{,}\ \lab_n\ \kw{:}\
\ty_n\ \kw{\kwrb}$ ($n \geq 0$), possibly enclosed in parentheses, and gives
the sequence of its labels in the order written:
\[
  \smlabs(\ty) = \lab_1 \cdots \lab_n .
\]
It is undefined on every other type expression, an extensible record type
expression or a type constructor among them.

The type names of a basis are intended to include those of its interfaces.
The following property can be demonstrated.

\begin{property}
Let $B$ be a basis with $\tynames(\ME \of B) \subseteq T \of B$, and let
$\jdg{B}{\modm}{(T', \ME', \PE'), E}$ be inferred. Then $T' \cap (T \of B) =
\emptyset$, and $\tynames E$, $\tynames \ME'$ and $\tynames \PE'$ are subsets
of $(T \of B) \cup T'$.

Moreover, in every instance of rule~\ruleref{r:sm-topdec-types}, $\TE$ is
determined by $C$ and the phrase, up to the choice of the type names in
$\Ran\TE$ that are not in $T \of C$.
\end{property}

\rulegroup{Modules}{\jdg{B}{\modm}{B', E}}

\defnrule[r:sm-module]{\jdg{B, \pi_0}{\imports_0\ \imports}{\pi} \premsep C = (T \of B,\ \emptyset,\ \smenv(\ME \of B, \pi)) \\ \jdg{C}{\topdec}{E, \PE} \premsep \jdg{E}{\header}{(\modid, I)} \\ \modid \notin \Dom(\ME \of B) \premsep T' = \tynames E \setminus (T \of B)}{\jdg{B}{\header\ \imports\ \topdec}{(T', \{\modid \mapsto I\}, \PE), E}}

\begin{rulecomments}
\rcomment{r:sm-module} The default imports are read before the imports of the
module, so an import written in the module may replace a type that a default
import exposes. % A28
The context $C$ holds only what the imports provide. The declarations of the
module are added by rule~\ruleref{r:sm-topdec-seq} to the unqualified
environment and never under a qualifier, so a module cannot refer to its own
names by qualification: in a module \kw{M}, \kw{M.x} refers to an import with
the qualifier \kw{M}, if there is one. % A38
The header is elaborated last, in the environment $E$ of the declarations of
the module alone; hence a module never exposes what it imports. The set $T'$
holds every type name that the module generates, private types included,
since an exposed value may have a type that mentions a private type. % A29
An import names a module of $\ME \of B$, and $\modid \notin \Dom(\ME \of B)$,
so no module imports itself and no cycle of imports elaborates;
Chapter~\ref{ch:programs} elaborates the modules of a program in an order in
which each follows the modules it imports. Nothing here constrains a value
named \kw{main}: only the \kw{main} of the main module of a program is
constrained, in Chapter~\ref{ch:programs}. % A32
\end{rulecomments}

\rulegroup{Headers}{\jdg{E}{\header}{(\modid, I)}}

\defnrule[r:sm-header]{I = \Expose(E, \explist)}{\jdg{E}{\opt{\kw{port}}\ \kw{module}\ \modid\ \kw{exposing}\ \explist}{(\modid, I)}}

\begin{rulecomments}
\rcomment{r:sm-header} Whether a module is a port module matters only to the
syntactic restrictions, which allow port declarations in it and require at
least one. The exposing list of a header elaborates by the operation $\Expose$
of Section~\ref{sec:static-modules-interfaces}, which the syntactic
restrictions make defined; the rules of the groups Exposing Lists and Exposing
Items below are for the exposing lists of imports.
\end{rulecomments}

\rulegroup{Imports}{\jdg{B, \pi}{\imports}{\pi'}}

In rule~\ruleref{r:sm-imports-import}, $\modid''$ is $\modid'$ if the first
option is present and $\modid$ if it is absent.

\defnrule[r:sm-imports-import]{(\ME \of B)(\modid) = I \premsep \pi_1 = \pi \smmerge_{\modid} (\emptyset,\ \modid''.\Dom I) \\ \optt{\jdg{\modid, I}{\explist}{(N, N')}} \premsep \pi_2 = \pi_1\ \optt{{}\smmerge_{\modid} (N, N')} \\ \jdg{B, \pi_2}{\imports}{\pi'}}{\jdg{B, \pi}{\kw{import}\ \modid\ \opt{\kw{as}\ \modid'}\ \optt{\kw{exposing}\ \explist}\ \imports}{\pi'}}

\defnrule[r:sm-imports-empty]{}{\jdg{B, \pi}{\ }{\pi}}

\begin{rulecomments}
\rcomment{r:sm-imports-import} Every identifier that the module exposes becomes
available under the qualifier $\modid''$, and the identifiers listed by the
exposing list become available unqualified. With \kw{as}, the qualifier is
$\modid'$ instead of $\modid$: after \kw{import Dict as D} the identifier
\kw{Dict.empty} is not available through this import. Several imports may
share a qualifier; their identifiers are joined under it, and an identifier
exposed under one qualifier by two modules is ambiguous. Since the imports are
read from left to right, a later import replaces what an earlier one provides
for the identifiers in its $N$. A module that is not in $\ME \of B$ cannot be
imported.
\end{rulecomments}

\rulegroup{Exposing Lists}{\jdg{\modid, I}{\explist}{(N, N')}}

\defnrule[r:sm-explist-all]{}{\jdg{\modid, I}{\kw{(..)}}{(\emptyset, \Dom I)}}

\defnrule[r:sm-explist-items]{\jdg{\modid, I}{\expitem_i}{(N_i, N'_i)} \quad (1 \le i \le n)}{\jdg{\modid, I}{\kw{(}\expitem_1\kw{,}\ \cdots\kw{,}\ \expitem_n\kw{)}}{(N_1 \cup \cdots \cup N_n,\ N'_1 \cup \cdots \cup N'_n)}}

\begin{rulecomments}
\rcomment{r:sm-explist-all} Everything is joined and nothing replaced, so two
imports with \kw{exposing (..)} of modules that both expose \kw{x} make
\kw{x} ambiguous. The imports elaborate nevertheless, and only a use of
\kw{x} fails to elaborate. % A28, A38
The operators of $\smops(\modid)$ are also exposed, in the sense of
Section~\ref{sec:syntax-modules-infix}, but no operator is bound.
\end{rulecomments}

\rulegroup{Exposing Items}{\jdg{\modid, I}{\expitem}{(N, N')}}

\defnrule[r:sm-item-var]{\var \in \Dom(\VE \of I)}{\jdg{\modid, I}{\var}{(\emptyset, \{\var\})}}

\defnrule[r:sm-item-op]{\opv \in \smops(\modid)}{\jdg{\modid, I}{\kw{(}\opv\kw{)}}{(\emptyset, \emptyset)}}

\defnrule[r:sm-item-type]{I(\tycon) = (\theta, \VE) \premsep N' = \{\con \,;\, \VE(\con) = (\sigma, \isv) \text{ for some } \sigma\}}{\jdg{\modid, I}{\tycon}{(\{\tycon\}, N')}}

\defnrule[r:sm-item-opentype]{I(\tycon) = (\theta, \VE) \premsep \VE(\con) = (\sigma, \isc) \text{ for some } \con, \sigma}{\jdg{\modid, I}{\tycon\,\kw{(..)}}{(\{\tycon\},\ \Dom\VE \cap \Dom(\VE \of I))}}

\begin{rulecomments}
\rcomment{r:sm-item-var} A value is joined with what earlier imports provide,
so a value exposed by name by two modules is ambiguous. % A28

\rcomment{r:sm-item-op} The rule only checks that the module declares the
operator, and the item binds nothing: the operators exposed to a module are
those that may be used in it in infix position
(Section~\ref{sec:syntax-modules-infix}), and the meaning of each is given by
$\OTstat$ (rule~\ruleref{r:s-atexp-op}), whatever the module imports.

\rcomment{r:sm-item-type} A type listed by name replaces any type of that name
provided by an earlier import: after \kw{import Foo exposing (Maybe)} the
type constructor \kw{Maybe} denotes the type of \kw{Foo}, not that of the
default import. % A28
For a record alias, $N'$ holds its constructor, which is joined; for a custom
type it is empty, so the constructors of a type are not imported by this item.
A constructor cannot be imported by name: in \kw{import Maybe exposing (Just)},
$I(\kw{Just})$ is undefined, and no rule applies.

\rcomment{r:sm-item-opentype} The side condition holds only for a custom type,
so \kw{Alias(..)} does not elaborate for an alias. For a type that its module
exposes closed, no constructor is in $\Dom(\VE \of I)$, so the item
elaborates, imports the type and imports no constructors: after
\kw{import A exposing (T(..))}, where \kw{A} exposes \kw{T} closed, the
constructors of \kw{T} are not available.
\end{rulecomments}

\rulegroup{Top-level Declarations}{\jdg{C}{\topdec}{E, \PE}}

\defnrule[r:sm-topdec-value]{\jdg{C + \VE}{\valbind}{\VE} \premsep \VE' = \Clos_C \VE}{\jdg{C}{\valbind}{\VE'\ \text{in}\ \Env,\ \emptymap}}

\defnrule[r:sm-topdec-types]{\opt{\jdg{C \oplus \TE}{\datbind}{\VE', \TE'} \premsep \forall (t, \VE'') \in \Ran\TE',\ t \notin T \of C} \\ \jdg{C \oplus \TE}{\typbind_i}{\VE_i, \TE_i} \quad (1 \le i \le l) \\ \TE = \opt{\TE' +{}} \TE_1 + \cdots + \TE_l \premsep \VE = \opt{\VE' +{}} \VE_1 + \cdots + \VE_l \\ \tyvars \TE = \emptyset}{C \vdash \opt{\kw{type}\ \datbind\ \kw{;}}\ \kw{type alias}\ \typbind_1\ \kw{;}\ \cdots \\ \qquad \kw{;}\ \kw{type alias}\ \typbind_l \Rightarrow (\emptymap, \TE, \VE),\ \emptymap}

\defnrule[r:sm-topdec-portout]{U \text{ is new for } \ty \text{ in } C \premsep \jdg{C + U}{\ty}{\tau} \\ \tau = \tau' \rightarrow \alpha\,\sbtCmd \premsep \Port(\tau')}{\jdg{C}{\kw{port}\ \var\ \kw{:}\ \ty}{\{\var \mapsto (\Clos_C \tau, \isv)\}\ \text{in}\ \Env,\ \{\var \mapsto (\smdirout, \tau')\}}}

\defnrule[r:sm-topdec-portin]{U \text{ is new for } \ty \text{ in } C \premsep \jdg{C + U}{\ty}{\tau} \\ \tau = (\tau' \rightarrow \alpha) \rightarrow \alpha\,\sbtSub \premsep \Port(\tau')}{\jdg{C}{\kw{port}\ \var\ \kw{:}\ \ty}{\{\var \mapsto (\Clos_C \tau, \isv)\}\ \text{in}\ \Env,\ \{\var \mapsto (\smdirin, \tau')\}}}

\defnrule[r:sm-topdec-seq]{\jdg{C}{\topdec_1}{E_1, \PE_1} \premsep \jdg{C \oplus E_1}{\topdec_2}{E_2, \PE_2} \\ \text{the separator is not one of a type-declaration phrase}}{\jdg{C}{\topdec_1\ \kw{;}\ \topdec_2}{E_1 + E_2,\ \PE_1 + \PE_2}}

\begin{rulecomments}
\rcomment{r:sm-topdec-value} The rule is that for declarations in the Core
(rule~\ruleref{r:s-dec-group}): a group of top-level values is recursive, an
unannotated member is monomorphic within its group, and the closure makes every
top-level value polymorphic in its other uses. % A24
At the top level the restriction on recursive values is syntactic and weaker
than in a \kw{let} (Section~\ref{sec:syntax-modules-restrictions}), so
\kw{f = \textbackslash n -> f n} elaborates here. % A26

\rcomment{r:sm-topdec-types} The phrase consists of the custom types of a
module, as one group, followed by its aliases, as Appendix~A arranges them
(Section~\ref{sec:appendix-derived-forms-modules}); here $l \geq 0$, and $l
\geq 1$ when the option is absent, in which case the phrase begins with the
first alias, and when $l = 0$ it is $\kw{type}\ \datbind$. Since $\TE$ occurs
on both sides, every type constructor of the phrase is in scope in all of it: a
constructor may take an argument whose type is an alias declared in the
phrase, and an alias may mention a custom type of the phrase or another alias.
The equation has a solution determined by $C$ and the phrase up to the choice
of new type names (Section~\ref{sec:static-modules-rules}), because the type
function of a custom type is its type name alone, and no alias refers to itself
through aliases. The side condition on $T \of C$ makes the type names of the
custom types new. The side condition $\tyvars\TE = \emptyset$ keeps the
equation from justifying itself: the closures of rules~\ruleref{r:sm-datbind}
and~\ruleref{r:sm-typbind-record} are taken in $C \oplus \TE$, so without it an
environment whose constructors keep a free variable, as \kw{K : a -> T a} for
\kw{type T a = K a}, would be a solution too; with it, the constructors are
closed and $\TE$ is determined as the property above says. % erratum, 2026-11-12
Rule~\ruleref{r:sm-topdec-seq} does not split such a
phrase (its side condition), so the type declarations of a module are
elaborated together, and their order does not matter: a custom type may
mention an alias that follows it. A local type spelt as an imported one is
a different type, and the unqualified type constructor denotes the local one,
in the declarations before it as well as after it: with \kw{type T = T String}
followed by \kw{type alias String = Int}, the constructor \kw{T} takes an
\kw{Int}. % A98

\rcomment{r:sm-topdec-portout} The variables of the annotation are made
explicit, as for an annotated value. The argument of $\sbtCmd$ must be a type
variable, so \kw{port send : String -> Cmd Msg} does not elaborate, nor does
\kw{port send : Cmd msg}, which has no payload, nor \kw{port send : Int -> Int
-> Cmd msg}. Since $\Port(\tau')$ excludes type variables, $\alpha$ is the only
variable of $\tau$, and the closure generalises it. The type $\sbtCmd$ is that
of the initial basis; a type of a program spelt \kw{Cmd} is a different type.
A port is a value of status $\isv$, so it is used and exposed like any other
value; $\PE$ records its direction and payload type, by which
Chapter~\ref{ch:effects} gives it its meaning, and Chapter~\ref{ch:programs}
requires the names of the ports of a program to be distinct. % A84
Thus \kw{port send : Char -> Cmd msg} does not elaborate, while \kw{port
send : List \kwlb\ id : Int, name : String \kwrb\ -> Cmd msg} does. % D11

\rcomment{r:sm-topdec-portin} The variable $\alpha$ occurs in both places, so
\kw{port recv : (Int -> a) -> Sub b} does not elaborate. As for an outgoing
port, \kw{port recv : (Json.Decode.Value -> msg) -> Sub msg} elaborates, and
\kw{port recv : (Result String Int -> msg) -> Sub msg} does not.

\rcomment{r:sm-topdec-seq} Each declaration is elaborated in a context that
contains the declarations before it. A type-declaration phrase, whose parts are
separated by \kw{;} too, is not split: rule~\ruleref{r:sm-topdec-types} elaborates
it as a whole, so that a custom type sees the aliases that follow it. % A98 Since $E_1$ modifies the environment of
$C$, a declaration of the module replaces an imported binding of the same
identifier, or makes an ambiguous one available. The syntactic restrictions
ensure that $\Dom E_1 \cap \Dom E_2 = \emptyset$ and $\Dom\PE_1 \cap \Dom\PE_2
= \emptyset$.
\end{rulecomments}

\rulegroup{Custom-type Bindings}{\jdg{C}{\datbind}{\VE, \TE}}

\defnrule[r:sm-datbind]{\tyvarseq = \tyrowvar^{(k)} \premsep U = \{\tyrowvar_1, \cdots, \tyrowvar_k\} \\ \jdg{C + U}{\tyvar_i}{\tau_i} \quad (1 \le i \le k) \\ \jdg{C + U, \tau^{(k)}t}{\conbind}{\VE} \premsep t \in \TyName^{(k)} \\ \opt{\jdg{C}{\datbind'}{\VE', \TE'} \premsep \forall (t', \VE'') \in \Ran\TE',\ t \neq t'}}{C \vdash \tycon\ \tyvarseq\ \kw{=}\ \conbind\ \opt{\kw{and}\ \datbind'} \Rightarrow \\ \qquad \Clos_C \VE\ \opt{{}+ \VE'},\ \{\tycon \mapsto (t, \Clos_C \VE)\}\ \opt{{}+ \TE'}}

\begin{rulecomments}
\rcomment{r:sm-datbind} The parameters are made explicit while the
constructors elaborate, and are generalised by the closure, since $U \of C =
\emptyset$ at the top level. A parameter need not occur in $\conbind$, so
phantom types are allowed. A parameter used as the extension variable of a
record type is a row variable; $\tau_i$ is then the open record type with no
fields of its own, and an argument of $t$ in that position is a record type.
The side condition $t \neq t'$ gives the custom types of one group distinct
type names; rule~\ruleref{r:sm-topdec-types} makes them new. The syntactic
restrictions ensure that $\Dom\VE \cap \Dom\VE' = \emptyset$ and $\tycon \notin
\Dom\TE'$.
\end{rulecomments}

\rulegroup{Constructor Bindings}{\jdg{C, \tau}{\conbind}{\VE}}

\defnrule[r:sm-conbind]{\jdg{C}{\ty_i}{\tau_i} \quad (1 \le i \le n) \premsep \opt{\jdg{C, \tau}{\conbind'}{\VE}}}{C, \tau \vdash \con\ \ty_1 \cdots \ty_n\ \opt{\kw{|}\ \conbind'} \Rightarrow \\ \qquad \{\con \mapsto (\tau_1 \rightarrow \cdots \rightarrow \tau_n \rightarrow \tau, \isc)\}\ \opt{{}+ \VE}}

\begin{rulecomments}
\rcomment{r:sm-conbind} When $n = 0$ the type of $\con$ is $\tau$, so a
constructor without arguments is a value of the custom type. The syntactic
restrictions ensure $\con \notin \Dom\VE$.
\end{rulecomments}

\rulegroup{Alias Bindings}{\jdg{C}{\typbind}{\VE, \TE}}

\defnrule[r:sm-typbind]{\tyvarseq = \tyrowvar^{(k)} \premsep \jdg{C + \{\tyrowvar_1, \cdots, \tyrowvar_k\}}{\ty}{\tau} \premsep \ty \notin \Dom\smlabs}{\jdg{C}{\tycon\ \tyvarseq\ \kw{=}\ \ty}{\emptymap,\ \{\tycon \mapsto (\Lambda\tyrowvar^{(k)}.\tau, \emptymap)\}}}

\defnrule[r:sm-typbind-record]{\tyvarseq = \tyrowvar^{(k)} \premsep \jdg{C + \{\tyrowvar_1, \cdots, \tyrowvar_k\}}{\ty}{\tau} \\ \smlabs(\ty) = \lab_1 \cdots \lab_n \premsep \tau = \rec{\varrho} \premsep \varrho = \{\lab_1 \mapsto \tau_1, \cdots, \lab_n \mapsto \tau_n\} \\ \con = \tycon \premsep \VE = \Clos_C \{\con \mapsto (\tau_1 \rightarrow \cdots \rightarrow \tau_n \rightarrow \tau, \isv)\}}{\jdg{C}{\tycon\ \tyvarseq\ \kw{=}\ \ty}{\VE,\ \{\tycon \mapsto (\Lambda\tyrowvar^{(k)}.\tau, \VE)\}}}

\begin{rulecomments}
\rcomment{r:sm-typbind} An alias generates no type name: it denotes a type
function, and a type expression that applies it denotes the expansion
(rule~\ruleref{r:s-ty-con}). Thus after \kw{type alias Id = Int}, the types
\kw{Id} and \kw{Int} are the same type. By the syntactic restrictions the
parameters are exactly the type variables of $\ty$, so the type function has no
unused parameter. Rules~\ruleref{r:sm-typbind}
and~\ruleref{r:sm-typbind-record} both cover the production; which applies is
decided by the form of $\ty$. An alias of an extensible record type, or of
another alias, defines no constructor, even when that alias is a record alias.

\rcomment{r:sm-typbind-record} A record alias also defines a constructor,
spelt as the alias, whose arguments are the fields in the order written: after
\kw{type alias P = \kwlb\ y : Int, x : Int \kwrb}, the expression \kw{P 1 2}
denotes the record \kw{\kwlb\ y = 1, x = 2 \kwrb}. Its status is $\isv$, so it
may be applied but may not occur in a pattern (rule~\ruleref{r:s-pat-con}).
It is exposed and imported together with the alias
(rule~\ruleref{r:sm-item-type}).
\end{rulecomments}

\section{Principal Derivations}\label{sec:static-modules-principal}

Instances and principal derivations are defined for a module as for a Core
phrase (Section~\ref{sec:static-core-principal}), with one addition: a module
generates type names, which a derivation chooses freely outside $T \of B$
(rules~\ruleref{r:sm-datbind} and~\ruleref{r:sm-topdec-types}). Let $D$ and
$D'$ be derivations for a module in a basis $B$. Then $D'$ is an
\emph{instance} of $D$ when, for some bijection $\rho$ from the type names that
$D$ generates to those that $D'$ generates, $D'$ is an instance, in the sense of
Section~\ref{sec:static-core-principal}, of the derivation obtained from $D$ by
replacing each such type name $t$ by $\rho(t)$; the sets $T$ of corresponding
contexts are then equal. A derivation for a module in $B$ is
\emph{principal} when every derivation for the module in $B$ is an instance of
it. % A96

\begin{property}
Every module that elaborates in a basis $B$ has a principal derivation in $B$,
and any two of its principal derivations differ only by a renaming of type and
row variables and of the type names they generate.
\end{property}

A derivation for a module combines derivations for its top-level declarations,
whose contexts the rules determine from $B$ and the earlier declarations, up
to the type names chosen; so the property follows from that of the Core. The
proof is not given here. % A95

\section{Further Restrictions}\label{sec:static-modules-further}

The further restrictions of the Core (Section~\ref{sec:static-core-further})
apply to the expressions and patterns of every top-level declaration, judged,
as in the Core, on the principal derivation for the module
(Section~\ref{sec:static-modules-principal}). % A96
The following conditions concern a program rather than a module, and
Chapter~\ref{ch:programs} imposes them: the type of \kw{main} in the main
module, a \kw{main} in any other module being an ordinary value; % A32
the distinctness of the names of the ports of a program from one another and
from the names of the managers of the platform
(Chapter~\ref{ch:effects}); % A84
the absence of any reference of \kw{main} to itself, directly or through a
delayed occurrence; and the distinctness of the identifiers of the modules of
a program from those of the modules of the initial basis. % A89

Some implementations accept further programs, such as those with a further
type of flags for \kw{main}, or reject programs that Standard Elm admits, such
as those with \kw{Bytes} in a payload type, or those that exceed limits on
the sizes of records, constructors and functions. None of these acceptances
or limits is part of Standard Elm; each is recorded in
Appendix~\ref{ch:appendix-implementation}. % A33, A34, A92
